\documentclass[10pt,a4paper]{article}
%% ДВА ДВИЖКА, ОДИН ФАЙЛ. У нас листок собирает LuaLaTeX (только он умеет
%% microtype целиком), а у коллег, которым .tex отдают на просмотр, стоит
%% обычный TeX Live с pdfLaTeX и Computer Modern — тем самым набором, которым
%% листки собирались в прошлом году. Файл обязан собираться и там, и там
%% (решение владельца 04.09). Ветка pdfLaTeX в песочнице не проверяется:
%% кириллических пакетов для него здесь нет, она написана по эталонному PDF.
\usepackage{iftex}
\ifPDFTeX
  \usepackage[T2A]{fontenc}
  \usepackage[utf8]{inputenc}
  \usepackage[russian]{babel}
\else
\usepackage{fontspec}
\usepackage{polyglossia}
\setmainlanguage{russian}
\setotherlanguage{english}
%% Начертания названы явно: без этого fontspec подставляет вместо курсива
%% ЖИРНЫЙ курсив, и все \textit в листке выходят жирными (проверено рендером).
%% Source Serif 4 — тот же шрифт, что выбран в html-конспектах проекта
%% (--serif в view.html). Файлы лежат рядом, в shrifty/, а не берутся
%% из системы: сборка должна давать один и тот же PDF в любой среде.
\setmainfont{SourceSerifPro}[
  Path           = shrifty/,
  Extension      = .ttf,
  UprightFont    = *_400Regular,
  ItalicFont     = *_400Regular_Italic,
  BoldFont       = *_700Bold,
  BoldItalicFont = *_700Bold_Italic]
\fi
\usepackage[left=11mm,right=13mm,top=18mm,bottom=14mm]{geometry}
\usepackage{amsmath,amssymb}
%% microtype с настройками из эталона владельца (_studio/konvejer/07-verstka/
%% etalon-tex/1.base-packages.sty): выступающая пунктуация и растяжение глифов
%% дают строке подобраться, и висящее на следующей строке слово втягивается
%% обратно — без этого приходится подгонять полосу руками.
%% Развилка по компилятору, а не выбор одного из двух. Целевой — LuaLaTeX:
%% на нём собран «Путеводитель Тьюринга» и только он умеет microtype целиком.
%% Но в песочнице, где собирается превью, luatex сломан (нет luatexbase.sty,
%% не грузятся шрифты Latin Modern, поставить неоткуда), и там работает
%% XeLaTeX — а он expansion и tracking не поддерживает и падает на них.
%% Так один и тот же .tex собирается в обеих средах, и в целевой — полностью.
\ifLuaTeX
  \usepackage[protrusion=true,expansion=true,tracking=true,
              stretch=50,shrink=30,final]{microtype}
\else
  \usepackage[protrusion=true,final]{microtype}
\fi
\emergencystretch=1em
\linespread{0.98}
\frenchspacing
\setlength{\parindent}{0pt}
\pagestyle{empty}

%% Окно для отметки перед задачей — как в листках 179 школы: ученик отмечает
%% сданное сам. Ширина подобрана по образцу 06_Графы.pdf.
%% Окно для плюса. Размеры из архива 179 (\shag=1.7cm, высота 0.5\baselineskip),
%% ширина ужата: там окно на весь плюс с подписью, здесь хватает плюса.
\newlength{\shirinaokna}\setlength{\shirinaokna}{12mm}
\newlength{\vysotaokna}\setlength{\vysotaokna}{0.7\baselineskip}
\newcommand{\okno}{\framebox[\shirinaokna]{\rule{0pt}{\vysotaokna}}}
\newlength{\shirinanomera}\setlength{\shirinanomera}{9mm}
\newlength{\zazor}\setlength{\zazor}{0.6em}

%% Абзац задачи. Висячего отступа НЕТ намеренно: в листках 179 продолжение
%% строки начинается у левого поля, а не под текстом первой строки, и весь
%% текст листка стоит на одном левом крае. Окно и номер — обычный текст
%% в начале первой строки, а не колонка на поле. Висячий вариант пробовался
%% 03.09 и отвергнут владельцем: «почему „доползти до любой другой“ должно
%% начинаться там же, где „из спичек“».
\newcommand{\blok}[1]{\par\noindent #1\par}

%% Номер в боксе постоянной ширины: без него однозначные и двузначные номера
%% дают колонку разной ширины, и текст задач стоит на двух разных вертикалях.
\newcommand{\nomer}[2]{\makebox[\shirinanomera][l]{%
  \fbox{\rule{0pt}{\vysotaokna}\textbf{#1}$^{#2}$}}}
%% Термин при первом употреблении. Одна команда на весь листок: выделение
%% значит «у слова специальный смысл», а не «мне тут захотелось курсива».
\newcommand{\term}[1]{\textit{#1}}
%% Пустое место вместо окна — там, где плюс ставится не за задачу, а за пункты.
\newcommand{\pusto}{\hspace*{\shirinaokna}}
%% Знак письменной сдачи. Стандартный теховский кинжал, а не картинка пера:
%% pifont-перо в кегле листка не читается вовсе (проверено рендером 02.09).
%% Знак письменной сдачи — надстрочный, чтобы не съедал строку и не лез
%% в текст условия.
\newcommand{\pero}{\textsuperscript{\dag}\;}
\newcommand{\zvezdy}{%
  \par\medskip
  \centerline{\rule{0.3\linewidth}{0.4pt}\quad $*$\ \ $*$\ \ $*$\quad
              \rule{0.3\linewidth}{0.4pt}}%
  \par\medskip}

\begin{document}
\centerline{\textit{Школа № 179. КЛ. 2026/2027 уч.год}}
\par\medskip\centerline{\rule{0.62\linewidth}{0.4pt}}\par\smallskip
\centerline{\large\textbf{1$\alpha$. \textsc{Деревья}}}\par\smallskip
\blok{\term{Деревом} называют связный граф без циклов.}\par\vspace{7pt}
\begin{samepage}
\blok{\nomer{1}{\circ}\hspace{\zazor}\okno\hspace{\zazor}Два графа \term{изоморфны}, если вершины одного можно занумеровать вершинами другого так, что рёбра перейдут в рёбра. Сколько существует попарно неизоморфных деревьев на шести вершинах? Нарисуйте их.}
\end{samepage}\par\vspace{7pt}
\begin{samepage}
\blok{\nomer{2}{\dag}\hspace{\zazor}\okno\hspace{\zazor}Вершину степени 1 называют \term{висячей}, или \term{листом}. Докажите, что в дереве, у которого больше одной вершины, есть висячая вершина, и что таких вершин хотя бы две.}
\end{samepage}\par\vspace{7pt}
\begin{samepage}
\blok{\nomer{3}{\dag}\hspace{\zazor}\okno\hspace{\zazor}Докажите, что в дереве с $n$ вершинами ровно $n-1$ ребро.}
\end{samepage}\par\vspace{7pt}
\begin{samepage}
\blok{\nomer{4}{\circ}\hspace{\zazor}\okno\hspace{\zazor}Сколько рёбер в \term{лесу} — графе без циклов — с $n$ вершинами и $k$ компонентами связности?}
\end{samepage}\par\vspace{7pt}
\begin{samepage}
\blok{\nomer{5}{\dag}\hspace{\zazor}\okno\hspace{\zazor}Докажите, что из любого связного графа можно выкинуть часть рёбер так, чтобы получилось дерево с теми же вершинами. Такое дерево называют \term{остовом}, или \term{остовным деревом}, графа.}
\end{samepage}\par\vspace{7pt}
\begin{samepage}
\blok{\nomer{6}{\circ}\hspace{\zazor}\okno\hspace{\zazor}В графе $n$ вершин и $m$ рёбер. Докажите, что компонент связности в нём не меньше $n-m$.}
\end{samepage}\par\vspace{7pt}
\begin{samepage}
\blok{\nomer{7}{\dag}\hspace{\zazor}\okno\hspace{\zazor}Пусть граф связен и имеет $n$ вершин. Докажите, что следующие свойства равносильны:}
\vspace{3pt}
\blok{\hspace{\zazor}\textbf{\large$\bullet$}\ рёбер ровно $n-1$;}
\blok{\hspace{\zazor}\textbf{\large$\bullet$}\ при удалении любого ребра граф перестаёт быть связным;}
\blok{\hspace{\zazor}\textbf{\large$\bullet$}\ в графе нет циклов;}
\blok{\hspace{\zazor}\textbf{\large$\bullet$}\ любые две вершины соединены ровно одним путём.}
\blok{Таким образом, любое из этих свойств можно взять за определение дерева.}
\end{samepage}\par\vspace{7pt}
\zvezdy
\begin{samepage}
\blok{\nomer{8}{\circ}\hspace{\zazor}\okno\hspace{\zazor}Даны $n$ различных по весу камней и весы, позволяющие сравнить любые два из них. Какое наименьшее число взвешиваний нужно, чтобы наверняка найти самый тяжёлый камень?}
\end{samepage}\par\vspace{7pt}
\begin{samepage}
\blok{\nomer{9}{}\hspace{\zazor}\okno\hspace{\zazor}Сколько рёбер может быть в несвязном графе с $n$ вершинами?}
\end{samepage}\par\vspace{7pt}
\begin{samepage}
\blok{\nomer{10}{\circ}\hspace{\zazor}\okno\hspace{\zazor}Докажите, что из любого связного графа, в котором больше одной вершины, можно удалить вершину вместе со всеми выходящими из неё рёбрами так, что граф останется связным.}
\end{samepage}\par\vspace{7pt}
\begin{samepage}
\blok{\nomer{11}{}\hspace{\zazor}\okno\hspace{\zazor}В связном графе 10 вершин и 19 рёбер. Докажите, что в нём найдётся цикл, после уничтожения всех рёбер которого граф не потеряет связность.}
\end{samepage}\par\vspace{7pt}
\begin{samepage}
\blok{\nomer{12}{}\hspace{\zazor}\okno\hspace{\zazor}В стране 100 городов, некоторые из которых соединены авиалиниями. Известно, что из любого города можно долететь до любого другого, возможно с пересадками. Докажите, что можно побывать в каждом городе, совершив не более 198 перелётов.}
\end{samepage}\par\vspace{7pt}
\begin{samepage}
\blok{\nomer{13}{}\hspace{\zazor}\okno\hspace{\zazor}Дано дерево с $n$ вершинами. Сколькими способами можно раскрасить его вершины в $k$ цветов так, чтобы концы любого ребра были разного цвета?}
\end{samepage}\par\vspace{7pt}
\begin{samepage}
\blok{\nomer{14}{}\hspace{\zazor}Выпуклый $n$-угольник разрезали на треугольники непересекающимися диагоналями.}
\blok{\okno\hspace{\zazor}\textbf{а)}$^{\circ}$ Докажите, что треугольников получилось на один больше, чем проведённых диагоналей.}
\blok{\okno\hspace{\zazor}\textbf{б)}$^{\circ}$ Докажите, что хотя бы у двух треугольников по две стороны лежат на границе многоугольника.}
\blok{\okno\hspace{\zazor}\textbf{в)}$^{\star}$ Сколькими способами можно так разрезать многоугольник? Вершины занумерованы: разрезания, отличающиеся поворотом, считаются различными.}
\end{samepage}\par\vspace{7pt}
\begin{samepage}
\blok{\nomer{15}{\star}\hspace{\zazor}\okno\hspace{\zazor}Хозяйка испекла пирог и знает, что придут либо ровно $p$, либо ровно $q$ гостей, причём $p$ и $q$ взаимно просты. На какое наименьшее число не обязательно равных кусков можно разрезать пирог, чтобы в любом случае разделить его между гостями поровну?}
\end{samepage}\par\vspace{7pt}
\end{document}